\documentclass[12pt]{article}
\usepackage[utf8]{inputenc}
\usepackage[T1]{fontenc}
\usepackage{hyperref}
\title{Umbu-Ungu Language Resources for K-12 Curriculum Development}
\author{Compiled by Claude Sonnet 5 (claude-sonnet-5) for Andr\'as Kornai}
\date{}
\begin{document}
\maketitle

\section{Introduction}
Umbu-Ungu (ISO 639-3 ubu; Glottocode umbu1258) is a Trans-New Guinea language of Papua New Guinea. It is often grouped with related varieties under ``Kaugel''. The sub-dialects named in the sources are Kala, No Penge and Andelale, and phonetic differences have led to separate written texts for them. Little published material was found. This dossier lists only what appeared in search results.

\section{Primary Educational Resources}
No primer, textbook or school curriculum for Umbu-Ungu turned up in the searches. Nothing indicates the language is taught in schools.

\section{Linguistic and Reference Materials}
\begin{itemize}
\item June Head, \emph{A Grammar of Umbu-Ungu} (2011). This is an unpublished SIL manuscript of about 249 pages, so it may be hard to obtain. It is listed in Glottolog: \url{https://glottolog.org/resource/reference/id/308940}
\item Grambank record for Umbu-Ungu, which cites the Head grammar: \url{https://grambank.clld.org/languages/umbu1258}
\item Numeral system entry (mBo-Ungu): \url{https://lingweb.eva.mpg.de/channumerals/Bo-Ung.htm}
\item KITLV thesaurus entry: \url{https://thesaurus-kitlv.library.leiden.edu/en/00006672.htm}
\end{itemize}
No published dictionary was found. A SIL linguist's message mentions a primary Umbu-Ungu dictionary, but the file was not seen. It may be available by contacting SIL.

\section{Religious and Community Resources}
The New Testament exists in two dialect editions: No Penge (1995) and Andelale (1998). Both are on eBible.org under Creative Commons terms (noncommercial, no derivatives).
\begin{itemize}
\item No Penge: \url{https://ebible.org/study/content/texts/ubu-nopenge/about.html}
\item No Penge (mirror): \url{https://cyber.bible/study/content/texts/ubu-nopenge/about.html}
\item Andelale: \url{https://ebible.org/study/content/texts/ubu-andelale/about.html}
\item Andelale (mirror): \url{https://cyber.bible/study/content/texts/ubu-andelale/about.html}
\end{itemize}

\section{Digital and Audio Resources}
\begin{itemize}
\item PARADISEC holds an Andelale wordlist recording with consultant Christopher Kinipi. Of 1,228 items, 729 words were recorded. \url{https://admin-catalog.paradisec.org.au/collections/OCSEAN2/items/UBU_20240802}
\end{itemize}

\section{Curriculum Development Context}
The materials above are enough to start a vocabulary and Scripture-based literacy component. A school curriculum would first need a dictionary and an agreed orthography across the dialects, and these were not found. The next step is to contact SIL Papua New Guinea about the grammar and dictionary, and to work with the community on the dialect question.
\end{document}